\chapter{Partial Trace and Feedback for Mealy Morphisms}
\label{ch:partial-trace-feedback-mealy}

\section{Overview}
\label{sec:partial-trace-feedback-overview}

The preceding chapters developed internal categories as monads in
\(\mathbf{Span}(\mathcal E)\), Mealy morphisms as stateful morphisms between
internal categories, and relative program categories as their induced execution
semantics.  The present chapter shows that this structure carries a canonical
partial trace.

The trace construction is not imposed as an additional operation.  It is
induced by the labelled input-discrete presentation of Mealy morphisms.  A
Mealy morphism
\[
\Phi:\mathbb A\rightsquigarrow\mathbb B
\]
is equivalently represented by a span of internal categories
\[
\begin{tikzcd}[column sep=large]
\mathbb A^{\mathrm{op}}
&
\mathbb P
  \arrow[l, "I"']
  \arrow[r, "\Omega"]
&
\mathbb B^{\mathrm{op}},
\end{tikzcd}
\]
whose left leg \(I\) is input-discrete.  Thus Mealy morphisms form a symmetric
monoidal subcategory of the span category of internal categories.

The ambient span category
\[
\mathbf{Span}(\mathbf{Cat}_{\mathrm{int}}(\mathcal E))
\]
has the standard total trace associated to spans.  The trace domain for Mealy
morphisms is forced by this embedding: a feedback composite is admissible
exactly when the ambient traced span still lies in the input-discrete
subcategory.  Equivalently, the hidden feedback equation has a unique internal
solution over each visible input and closed state.

The main construction is as follows.  Given a Mealy morphism
\[
\Phi:\mathbb A\times\mathbb U
\rightsquigarrow
\mathbb B\times\mathbb U,
\]
write its labelled input-discrete presentation as
\[
\begin{tikzcd}[column sep=large]
(\mathbb A\times\mathbb U)^{\mathrm{op}}
&
\mathbb P
  \arrow[l, "I"']
  \arrow[r, "\Omega"]
&
(\mathbb B\times\mathbb U)^{\mathrm{op}} .
\end{tikzcd}
\]
Using the canonical product-opposite identifications
\[
(\mathbb A\times\mathbb U)^{\mathrm{op}}
\cong
\mathbb A^{\mathrm{op}}\times\mathbb U^{\mathrm{op}},
\qquad
(\mathbb B\times\mathbb U)^{\mathrm{op}}
\cong
\mathbb B^{\mathrm{op}}\times\mathbb U^{\mathrm{op}},
\]
write
\[
I=(I_A,I_U),
\qquad
\Omega=(\Omega_B,\Omega_U).
\]
The ambient span trace closes the hidden \(\mathbb U\)-interface by forming
the equalizer
\[
\mathbb P^\circ
:=
\operatorname{Eq}(I_U,\Omega_U).
\]
The trace is defined precisely when the remaining visible input projection
\[
I_A:\mathbb P^\circ\to\mathbb A^{\mathrm{op}}
\]
is input-discrete.  In that case, the traced Mealy morphism is represented by
the span
\[
\begin{tikzcd}[column sep=large]
\mathbb A^{\mathrm{op}}
&
\mathbb P^\circ
  \arrow[l, "I_A"']
  \arrow[r, "\Omega_B"]
&
\mathbb B^{\mathrm{op}} .
\end{tikzcd}
\]

In component notation, if
\[
\Phi=(P,\delta,\omega):
\mathbb A\times\mathbb U
\rightsquigarrow
\mathbb B\times\mathbb U,
\]
the trace-class condition says that the internal feedback equation
\[
h=\omega_U((a,h),x)
\]
has a unique solution for every closed state \(x\) and every visible input
arrow \(a\) admissible at \(x\).  Thus trace-class morphisms are precisely
those Mealy morphisms for which the hidden feedback interface can be closed
while remaining deterministic in the Mealy sense.

The central slogan of the chapter is therefore:
\[
\boxed{
\text{partial trace for Mealy morphisms}
=
\text{ambient span trace restricted to input-discrete spans}.
}
\]
Equivalently, in component form:
\[
\boxed{
\text{trace-class}
=
\text{unique solvability of the Mealy feedback equation}.
}
\]

The chapter also proves the compatibility of staged and simultaneous feedback.
For a morphism with two hidden interfaces,
\[
\Phi:
\mathbb A\times\mathbb U\times\mathbb V
\rightsquigarrow
\mathbb B\times\mathbb U\times\mathbb V,
\]
staged feedback agrees with simultaneous feedback whenever the staged feedback
problems are trace-class.  This is the Mealy form of the vanishing law for the
ambient span trace.

\section{Standing hypotheses and span conventions}
\label{sec:partial-trace-standing-hypotheses}

Throughout this chapter, assume that \(\mathcal E\) has finite limits.  Then
the category
\[
\mathbf{Cat}_{\mathrm{int}}(\mathcal E)
\]
of internal categories and internal functors has finite limits, computed
componentwise.  In particular, products, pullbacks, and equalizers of internal
categories exist and are created by the object-of-objects and object-of-arrows
projections.

The partial-trace construction itself uses only these finite-limit hypotheses.
Whenever predicate-transformer, fixed-point, or polynomial-response structures
from previous chapters are invoked, the corresponding additional hypotheses
from those chapters are assumed.

We write
\[
\mathbf{Span}(\mathbf{Cat}_{\mathrm{int}}(\mathcal E))
\]
for the ordinary span category whose morphisms are isomorphism classes of spans
of internal categories.  Equivalently, one may work with chosen pullbacks and
read all span equalities up to the canonical apex isomorphisms.

An ordinary span in this category is displayed as
\[
\begin{tikzcd}[column sep=large]
\mathbb X
&
\mathbb R
  \arrow[l, "\ell"']
  \arrow[r, "r"]
&
\mathbb Y .
\end{tikzcd}
\]
Given spans
\[
\begin{tikzcd}[column sep=large]
\mathbb X
&
\mathbb R
  \arrow[l, "\ell"']
  \arrow[r, "r"]
&
\mathbb Y
\end{tikzcd}
\qquad
\text{and}
\qquad
\begin{tikzcd}[column sep=large]
\mathbb Y
&
\mathbb S
  \arrow[l, "s"']
  \arrow[r, "u"]
&
\mathbb Z,
\end{tikzcd}
\]
their composite is represented by
\[
\begin{tikzcd}[column sep=large]
\mathbb X
&
\mathbb R\times_{\mathbb Y}\mathbb S
  \arrow[l, "\ell\pi_{\mathbb R}"']
  \arrow[r, "u\pi_{\mathbb S}"]
&
\mathbb Z .
\end{tikzcd}
\]

All references to input-discreteness use the convention fixed in the preceding
chapters.  Thus an internal functor
\[
J:\mathbb X\to\mathbb A^{\mathrm{op}}
\]
is input-discrete when the square
\[
\begin{tikzcd}[column sep=large,row sep=large]
X_1
  \arrow[r,"J_1"]
  \arrow[d,"s_X"']
&
A_1
  \arrow[d,"t_A"]
\\
X_0
  \arrow[r,"J_0"']
&
A_0
\end{tikzcd}
\]
is a pullback.  This convention is adapted to right internal actions: an arrow
\((a,x)\) starts at a state \(x\) lying over \(t_A(a)\), and executes along
\(a\) to a state lying over \(s_A(a)\).

\section{The Mealy category inside spans of internal categories}
\label{sec:mealy-category-inside-spans}

We write
\[
\mathsf{Mly}(\mathcal E)
\]
for the ordinary category whose objects are internal categories in
\(\mathcal E\) and whose morphisms are Mealy morphisms, taken up to the
canonical isomorphisms of carrier spans generated by chosen pullbacks.
Composition and identities are those constructed in
Chapter~\ref{ch:spans-internal-categories-mealy}.

The tensor product on objects is product of internal categories.  If
\[
\Phi=(P,\delta_\Phi,\omega_\Phi):
\mathbb A\rightsquigarrow\mathbb B
\]
and
\[
\Psi=(Q,\delta_\Psi,\omega_\Psi):
\mathbb C\rightsquigarrow\mathbb D,
\]
then their tensor product
\[
\Phi\times\Psi:
\mathbb A\times\mathbb C
\rightsquigarrow
\mathbb B\times\mathbb D
\]
has carrier \(P\times Q\), transition map
\[
\delta_{\Phi\times\Psi}
=
\delta_\Phi\times\delta_\Psi,
\]
and output map
\[
\omega_{\Phi\times\Psi}
=
\omega_\Phi\times\omega_\Psi.
\]

\begin{proposition}
\label{prop:mly-symmetric-monoidal}
The category
\[
\mathsf{Mly}(\mathcal E)
\]
is symmetric monoidal under product of internal categories.
\end{proposition}

\begin{proof}
Products of internal categories are computed componentwise, and the product of
two Mealy morphisms is defined componentwise on carriers, transitions, and
outputs.  The Mealy unit and multiplication laws hold componentwise.  The
associativity, unit, and symmetry constraints are inherited from products of
internal categories and from the span calculus.  Hence
\(\mathsf{Mly}(\mathcal E)\) is symmetric monoidal.
\end{proof}

Recall that a Mealy morphism
\[
\Phi:\mathbb A\rightsquigarrow\mathbb B
\]
is equivalently a labelled input-discrete internal category
\[
\begin{tikzcd}[column sep=large]
&
\mathbb P
  \arrow[dl, "I"']
  \arrow[dr, "\Omega"]
\\
\mathbb A^{\mathrm{op}}
&&
\mathbb B^{\mathrm{op}},
\end{tikzcd}
\]
with \(I\) input-discrete.  Equivalently, it is a span of internal categories
\[
\begin{tikzcd}[column sep=large]
\mathbb A^{\mathrm{op}}
&
\mathbb P
  \arrow[l, "I"']
  \arrow[r, "\Omega"]
&
\mathbb B^{\mathrm{op}}
\end{tikzcd}
\]
whose left leg is input-discrete.

\begin{definition}[The span subcategory of Mealy morphisms]
\label{def:span-subcategory-mealy}
Let
\[
\mathcal S_{\mathrm{Mly}}
\subseteq
\mathbf{Span}(\mathbf{Cat}_{\mathrm{int}}(\mathcal E))
\]
be the subcategory whose objects are internal categories of the form
\(\mathbb A^{\mathrm{op}}\), and whose morphisms are spans
\[
\begin{tikzcd}[column sep=large]
\mathbb A^{\mathrm{op}}
&
\mathbb P
  \arrow[l, "I"']
  \arrow[r, "\Omega"]
&
\mathbb B^{\mathrm{op}}
\end{tikzcd}
\]
with \(I\) input-discrete, taken up to labelled isomorphism of apex internal
categories.
\end{definition}

\begin{theorem}[Mealy morphisms as input-discrete spans]
\label{thm:mealy-as-input-discrete-spans}
The labelled input-discrete construction gives a symmetric monoidal
equivalence of ordinary categories
\[
\mathsf{Mly}(\mathcal E)
\simeq
\mathcal S_{\mathrm{Mly}}.
\]
Composed with the inclusion
\[
\mathcal S_{\mathrm{Mly}}
\hookrightarrow
\mathbf{Span}(\mathbf{Cat}_{\mathrm{int}}(\mathcal E)),
\]
this gives a faithful symmetric monoidal span representation of Mealy
morphisms.
\end{theorem}

\begin{proof}
The equivalence on hom-sets is the labelled input-discrete characterization of
Mealy morphisms established in
Chapter~\ref{ch:relative-mealy-program-categories}.  The quotient by labelled
isomorphism of apex internal categories matches the quotient by canonical
carrier-span isomorphisms in \(\mathsf{Mly}(\mathcal E)\).

Identities correspond to identity input-discrete spans.  Composition
corresponds to pullback of labelled input-discrete spans.  Given spans
\[
\begin{tikzcd}[column sep=large]
\mathbb A^{\mathrm{op}}
&
\mathbb P
  \arrow[l, "I"']
  \arrow[r, "\Omega"]
&
\mathbb B^{\mathrm{op}}
\end{tikzcd}
\qquad
\text{and}
\qquad
\begin{tikzcd}[column sep=large]
\mathbb B^{\mathrm{op}}
&
\mathbb Q
  \arrow[l, "J"']
  \arrow[r, "\Theta"]
&
\mathbb C^{\mathrm{op}},
\end{tikzcd}
\]
their composite has apex
\[
\mathbb P\times_{\mathbb B^{\mathrm{op}}}\mathbb Q.
\]
The left leg of the composite is
\[
I\pi_{\mathbb P}:
\mathbb P\times_{\mathbb B^{\mathrm{op}}}\mathbb Q
\longrightarrow
\mathbb A^{\mathrm{op}}.
\]

Concretely, given an object \((p,q)\) of the pullback apex and an
\(\mathbb A\)-arrow \(a\) with \(t_A(a)=I_0(p)\), input-discreteness of \(I\)
gives a unique arrow of \(\mathbb P\) over \(a\) with source \(p\).  Its
\(\mathbb B^{\mathrm{op}}\)-label then determines, by input-discreteness of
\(J\), a unique compatible arrow of \(\mathbb Q\) with source \(q\).  Hence the
composite left leg is input-discrete.

The construction is compatible with products, since products of internal
categories are computed componentwise and products of pullback squares are
pullback squares.  Therefore the equivalence is symmetric monoidal.
\end{proof}

\section{The ambient total trace on spans of internal categories}
\label{sec:ambient-total-trace-spans}

The ambient category
\[
\mathbf{Span}(\mathbf{Cat}_{\mathrm{int}}(\mathcal E))
\]
has the standard total trace of spans with respect to the product monoidal
structure.

Let
\[
\begin{tikzcd}[column sep=large]
\mathbb X\times\mathbb U
&
\mathbb R
  \arrow[l, "\ell"']
  \arrow[r, "r"]
&
\mathbb Y\times\mathbb U
\end{tikzcd}
\]
be a span of internal categories.  Write
\[
\ell=(\ell_X,\ell_U),
\qquad
r=(r_Y,r_U).
\]
The trace over \(\mathbb U\) is the span
\[
\begin{tikzcd}[column sep=large]
\mathbb X
&
\mathbb R^\circ
  \arrow[l, "\ell_X\iota"']
  \arrow[r, "r_Y\iota"]
&
\mathbb Y,
\end{tikzcd}
\]
where
\[
\mathbb R^\circ
:=
\operatorname{Eq}(\ell_U,r_U),
\]
with inclusion
\[
\iota:\mathbb R^\circ\hookrightarrow\mathbb R.
\]
Equivalently,
\[
\mathbb R^\circ
=
\mathbb R\times_{\mathbb U\times\mathbb U}\mathbb U,
\]
using
\[
(\ell_U,r_U):\mathbb R\to\mathbb U\times\mathbb U
\]
and the diagonal
\[
\Delta_{\mathbb U}:\mathbb U\to\mathbb U\times\mathbb U.
\]

\begin{proposition}[Total trace of spans]
\label{prop:ambient-span-total-trace}
The operation above is the standard total trace on
\[
\mathbf{Span}(\mathbf{Cat}_{\mathrm{int}}(\mathcal E))
\]
with respect to the product monoidal structure.
\end{proposition}

\begin{proof}
The span category of any finitely complete category is compact closed under
the product monoidal structure.  The dual of an object \(\mathbb U\) is
\(\mathbb U\) itself.  The coevaluation and evaluation are represented by the
diagonal spans
\[
\mathbf 1
\longleftarrow
\mathbb U
\xrightarrow{\Delta_{\mathbb U}}
\mathbb U\times\mathbb U
\]
and
\[
\mathbb U\times\mathbb U
\xleftarrow{\Delta_{\mathbb U}}
\mathbb U
\longrightarrow
\mathbf 1.
\]
The trace of a span
\[
\mathbb X\times\mathbb U
\xleftarrow{\ell}
\mathbb R
\xrightarrow{r}
\mathbb Y\times\mathbb U
\]
is obtained by composing with these evaluation and coevaluation spans.  The
resulting apex is canonically isomorphic to
\[
\mathbb R\times_{\mathbb U\times\mathbb U}\mathbb U,
\]
where
\[
\mathbb R\to\mathbb U\times\mathbb U
\]
is \((\ell_U,r_U)\) and
\[
\mathbb U\to\mathbb U\times\mathbb U
\]
is the diagonal.  This is precisely
\[
\operatorname{Eq}(\ell_U,r_U).
\]
The trace axioms are the yanking, sliding, superposing, and vanishing
identities of the compact closed trace, and hold up to the canonical span apex
isomorphisms.
\end{proof}

\section{Trace class for Mealy morphisms}
\label{sec:trace-class-mealy-morphisms}

Let
\[
\Phi:
\mathbb A\times\mathbb U
\rightsquigarrow
\mathbb B\times\mathbb U
\]
be a Mealy morphism.  Present it as a labelled input-discrete span
\[
\begin{tikzcd}[column sep=large]
(\mathbb A\times\mathbb U)^{\mathrm{op}}
&
\mathbb P
  \arrow[l, "I"']
  \arrow[r, "\Omega"]
&
(\mathbb B\times\mathbb U)^{\mathrm{op}} .
\end{tikzcd}
\]
Using the canonical identifications
\[
(\mathbb A\times\mathbb U)^{\mathrm{op}}
\cong
\mathbb A^{\mathrm{op}}\times\mathbb U^{\mathrm{op}},
\qquad
(\mathbb B\times\mathbb U)^{\mathrm{op}}
\cong
\mathbb B^{\mathrm{op}}\times\mathbb U^{\mathrm{op}},
\]
write
\[
I=(I_A,I_U),
\qquad
\Omega=(\Omega_B,\Omega_U).
\]

\begin{definition}[Closed-loop category]
\label{def:closed-loop-category}
The \textbf{closed-loop category} of \(\Phi\) over \(\mathbb U\) is
\[
\mathbb P^\circ
:=
\operatorname{Eq}(I_U,\Omega_U).
\]
Thus \(\mathbb P^\circ\) is the internal subcategory of \(\mathbb P\) whose
hidden \(\mathbb U\)-input and hidden \(\mathbb U\)-output labels agree.
\end{definition}

Equivalently, \(\mathbb P^\circ\) is characterized by the pullback square
\[
\begin{tikzcd}[column sep=large,row sep=large]
\mathbb P^\circ
  \arrow[r, hook, "\iota"]
  \arrow[d]
&
\mathbb P
  \arrow[d, "{(I_U,\Omega_U)}"]
\\
\mathbb U^{\mathrm{op}}
  \arrow[r, "\Delta_{\mathbb U^{\mathrm{op}}}"']
&
\mathbb U^{\mathrm{op}}\times\mathbb U^{\mathrm{op}} .
\end{tikzcd}
\]

The ambient span trace of the labelled input-discrete span representing
\(\Phi\) is
\[
\begin{tikzcd}[column sep=large]
\mathbb A^{\mathrm{op}}
&
\mathbb P^\circ
  \arrow[l, "I_A\iota"']
  \arrow[r, "\Omega_B\iota"]
&
\mathbb B^{\mathrm{op}} .
\end{tikzcd}
\]
We suppress \(\iota\) from the notation when no confusion can arise.

Although the ambient span trace is formally taken over
\[
\mathbb U^{\mathrm{op}},
\]
we denote the resulting Mealy trace by
\[
\operatorname{Tr}^{\mathbb U},
\]
since \(\mathbb U\) is the hidden interface in the original Mealy morphism.

\begin{definition}[Trace-class Mealy morphism]
\label{def:trace-class-mealy}
The Mealy morphism
\[
\Phi:
\mathbb A\times\mathbb U
\rightsquigarrow
\mathbb B\times\mathbb U
\]
is \textbf{trace-class over \(\mathbb U\)} if
\[
I_A:\mathbb P^\circ\to\mathbb A^{\mathrm{op}}
\]
is input-discrete.
\end{definition}

\begin{definition}[Partial trace of a trace-class Mealy morphism]
\label{def:partial-trace-mealy}
If \(\Phi\) is trace-class over \(\mathbb U\), its \textbf{partial trace} is the
Mealy morphism
\[
\operatorname{Tr}^{\mathbb U}(\Phi):
\mathbb A\rightsquigarrow\mathbb B
\]
represented by the labelled input-discrete span
\[
\begin{tikzcd}[column sep=large]
\mathbb A^{\mathrm{op}}
&
\mathbb P^\circ
  \arrow[l, "I_A"']
  \arrow[r, "\Omega_B"]
&
\mathbb B^{\mathrm{op}} .
\end{tikzcd}
\]
Equivalently, this is the ambient span trace over
\(\mathbb U^{\mathrm{op}}\), transported through the opposite convention used
in the Mealy-span representation.
\end{definition}

\begin{proposition}[Trace class as closure under ambient trace]
\label{prop:trace-class-ambient-closure}
For a Mealy morphism
\[
\Phi:
\mathbb A\times\mathbb U
\rightsquigarrow
\mathbb B\times\mathbb U,
\]
the following are equivalent:
\begin{enumerate}
\item \(\Phi\) is trace-class over \(\mathbb U\);
\item the ambient span trace of the labelled input-discrete span associated to
\(\Phi\) lies again in the Mealy span subcategory
\[
\mathcal S_{\mathrm{Mly}}
\subseteq
\mathbf{Span}(\mathbf{Cat}_{\mathrm{int}}(\mathcal E)).
\]
\end{enumerate}
\end{proposition}

\begin{proof}
The ambient span trace is
\[
\begin{tikzcd}[column sep=large]
\mathbb A^{\mathrm{op}}
&
\mathbb P^\circ
  \arrow[l, "I_A"']
  \arrow[r, "\Omega_B"]
&
\mathbb B^{\mathrm{op}} .
\end{tikzcd}
\]
By definition of \(\mathcal S_{\mathrm{Mly}}\), this span lies in the Mealy
subcategory exactly when its left leg \(I_A\) is input-discrete.  This is
precisely the trace-class condition.
\end{proof}

\section{The internal feedback equation}
\label{sec:internal-feedback-equation}

We now unpack the trace-class condition in the original component notation of
Mealy morphisms.

Let
\[
\Phi=(P,\delta,\omega):
\mathbb A\times\mathbb U
\rightsquigarrow
\mathbb B\times\mathbb U.
\]
The carrier span has the form
\[
\begin{tikzcd}[column sep=large]
A_0\times U_0
&
P
  \arrow[l, "p"']
  \arrow[r, "q"]
&
B_0\times U_0 .
\end{tikzcd}
\]
Write
\[
p=(p_A,p_U),
\qquad
q=(q_B,q_U).
\]
The transition and output maps are defined on admissible triples
\[
(a,h,x)
\]
with
\[
t_A(a)=p_A(x),
\qquad
t_U(h)=p_U(x).
\]
The output component decomposes as
\[
\omega=(\omega_B,\omega_U).
\]

\begin{definition}[Closed state object]
\label{def:closed-state-object}
The \textbf{closed state object} of \(\Phi\) over \(\mathbb U\) is
\[
P^\circ:=\operatorname{Eq}(p_U,q_U).
\]
A generalized element of \(P^\circ\) is a state \(x\in P\) satisfying
\[
p_U(x)=q_U(x).
\]
\end{definition}

Equivalently, \(P^\circ\) is given by the pullback
\[
\begin{tikzcd}[column sep=large,row sep=large]
P^\circ
  \arrow[r, hook, "\iota_P"]
  \arrow[d]
&
P
  \arrow[d, "{(p_U,q_U)}"]
\\
U_0
  \arrow[r, "\Delta_{U_0}"']
&
U_0\times U_0 .
\end{tikzcd}
\]

Let
\[
D_\Phi
:=
A_1\times_{t_A,A_0,p_A}P^\circ.
\]
A generalized element of \(D_\Phi\) is an admissible visible input/closed-state
pair
\[
(a,x)
\]
with
\[
t_A(a)=p_A(x).
\]

Define
\[
H_\Phi
:=
D_\Phi\times_{p_U\pi_P,U_0,t_U}U_1.
\]
A generalized element of \(H_\Phi\) is a triple
\[
(a,x,h)
\]
such that
\[
t_A(a)=p_A(x),
\qquad
t_U(h)=p_U(x)=q_U(x).
\]

There are two maps
\[
H_\Phi\to U_1.
\]
The first is the projection
\[
\pi_U:H_\Phi\to U_1,
\qquad
(a,x,h)\longmapsto h.
\]
The second is induced by the hidden output component:
\[
\chi_\Phi:H_\Phi\to U_1,
\qquad
(a,x,h)\longmapsto\omega_U((a,h),x).
\]
This map is well-defined because \((a,h)\) is an admissible arrow of
\(\mathbb A\times\mathbb U\) at \(x\), and because
\[
t_U\omega_U((a,h),x)=q_U(x)=p_U(x)=t_U(h).
\]

\begin{definition}[Feedback witness object]
\label{def:feedback-witness-object}
The \textbf{feedback witness object} of \(\Phi\) over \(\mathbb U\) is the
equalizer
\[
W_\Phi
:=
\operatorname{Eq}(\pi_U,\chi_\Phi)
\hookrightarrow
H_\Phi.
\]
Thus a generalized element of \(W_\Phi\) is a triple
\[
(a,x,h)
\]
with
\[
t_A(a)=p_A(x),
\qquad
t_U(h)=p_U(x)=q_U(x),
\qquad
h=\omega_U((a,h),x).
\]
There is a canonical projection
\[
\pi_\Phi:W_\Phi\to D_\Phi,
\qquad
(a,x,h)\longmapsto(a,x).
\]
\end{definition}

\begin{theorem}[Internal trace-class criterion]
\label{thm:internal-trace-class-criterion}
For
\[
\Phi:
\mathbb A\times\mathbb U
\rightsquigarrow
\mathbb B\times\mathbb U,
\]
the following are equivalent:
\begin{enumerate}
\item \(\Phi\) is trace-class over \(\mathbb U\);
\item the functor
\[
I_A:\mathbb P^\circ\to\mathbb A^{\mathrm{op}}
\]
is input-discrete;
\item the projection
\[
\pi_\Phi:W_\Phi\to D_\Phi
\]
is an isomorphism.
\end{enumerate}
\end{theorem}

\begin{proof}
The equivalence between (1) and (2) is the definition of trace-class.

It remains to identify input-discreteness of
\[
I_A:\mathbb P^\circ\to\mathbb A^{\mathrm{op}}
\]
with the assertion that
\[
\pi_\Phi:W_\Phi\to D_\Phi
\]
is an isomorphism.

Since
\[
I:\mathbb P\to(\mathbb A\times\mathbb U)^{\mathrm{op}}
\]
is input-discrete, the square
\[
\begin{tikzcd}[column sep=large,row sep=large]
P_1
  \arrow[r, "I_1"]
  \arrow[d, "s_P"']
&
A_1\times U_1
  \arrow[d, "t_A\times t_U"]
\\
P_0
  \arrow[r, "p_A\times p_U"']
&
A_0\times U_0
\end{tikzcd}
\]
is a pullback.  Hence
\[
P_1
\cong
(A_1\times U_1)
\times_{t_A\times t_U,A_0\times U_0,p_A\times p_U}
P.
\]
Thus an arrow of \(\mathbb P\) is represented uniquely by a triple
\[
(a,h,x)
\]
with
\[
t_A(a)=p_A(x),
\qquad
t_U(h)=p_U(x).
\]
Its target in \(\mathbb P\) is
\[
\delta((a,h),x),
\]
and its output label in
\[
(\mathbb B\times\mathbb U)^{\mathrm{op}}
\]
is
\[
(\omega_B((a,h),x),\omega_U((a,h),x)).
\]

The closed-loop category
\[
\mathbb P^\circ=\operatorname{Eq}(I_U,\Omega_U)
\]
has object object
\[
P^\circ=\operatorname{Eq}(p_U,q_U).
\]
Its arrow object is obtained by equalizing the hidden input and hidden output
labels on arrows.  Therefore, under the displayed identification of \(P_1\),
an arrow of \(\mathbb P^\circ\) is precisely a triple
\[
(a,h,x)
\]
such that
\[
x\in P^\circ,
\qquad
t_A(a)=p_A(x),
\qquad
t_U(h)=p_U(x),
\qquad
h=\omega_U((a,h),x).
\]
Since \(x\in P^\circ\), the condition \(t_U(h)=p_U(x)\) is equivalently
\[
t_U(h)=p_U(x)=q_U(x).
\]
Hence
\[
P^\circ_1\cong W_\Phi.
\]

The defining input-discreteness square for
\[
I_A:\mathbb P^\circ\to\mathbb A^{\mathrm{op}}
\]
has comparison map
\[
P^\circ_1
\longrightarrow
A_1\times_{t_A,A_0,p_A}P^\circ
=
D_\Phi.
\]
Under the identification
\[
P^\circ_1\cong W_\Phi,
\]
this comparison map is exactly
\[
\pi_\Phi:W_\Phi\to D_\Phi.
\]
Therefore \(I_A\) is input-discrete if and only if
\(\pi_\Phi\) is an isomorphism.
\end{proof}

\begin{remark}[Trace class as unique feedback solvability]
\label{rem:trace-class-unique-feedback}
The ambient span trace keeps all hidden witnesses satisfying the feedback
equation
\[
h=\omega_U((a,h),x).
\]
The Mealy trace is defined exactly when this hidden-witness relation is
functional and total over visible admissible inputs and closed states.  Thus
trace class is the internal unique-solvability condition for the feedback
equation.
\end{remark}

When \(\Phi\) is trace-class, the isomorphism
\[
\pi_\Phi:W_\Phi\to D_\Phi
\]
defines an internal morphism
\[
h_\Phi:D_\Phi\to U_1
\]
by
\[
h_\Phi=\pi_U\circ\pi_\Phi^{-1}.
\]
Thus \(h_\Phi(a,x)\) is the unique hidden arrow satisfying
\[
h_\Phi(a,x)=\omega_U((a,h_\Phi(a,x)),x).
\]

\section{Component formula for the partial trace}
\label{sec:component-formula-partial-trace}

\begin{proposition}[Component formula for the partial trace]
\label{prop:component-formula-partial-trace}
Let
\[
\Phi=(P,\delta,\omega):
\mathbb A\times\mathbb U
\rightsquigarrow
\mathbb B\times\mathbb U
\]
be trace-class over \(\mathbb U\).  Then
\[
\operatorname{Tr}^{\mathbb U}(\Phi):
\mathbb A\rightsquigarrow\mathbb B
\]
has carrier span
\[
\begin{tikzcd}[column sep=large]
A_0
&
P^\circ
  \arrow[l, "p_A"']
  \arrow[r, "q_B"]
&
B_0,
\end{tikzcd}
\]
and its transition and output maps are
\[
\delta^{\operatorname{tr}}(a,x)
=
\delta((a,h_\Phi(a,x)),x),
\]
and
\[
\omega^{\operatorname{tr}}(a,x)
=
\omega_B((a,h_\Phi(a,x)),x).
\]
\end{proposition}

\begin{proof}
By definition, the traced Mealy morphism is represented by the input-discrete
span
\[
\begin{tikzcd}[column sep=large]
\mathbb A^{\mathrm{op}}
&
\mathbb P^\circ
  \arrow[l, "I_A"']
  \arrow[r, "\Omega_B"]
&
\mathbb B^{\mathrm{op}} .
\end{tikzcd}
\]
The object part of this span is
\[
\begin{tikzcd}[column sep=large]
A_0
&
P^\circ
  \arrow[l, "p_A"']
  \arrow[r, "q_B"]
&
B_0.
\end{tikzcd}
\]

Since \(I_A\) is input-discrete, for each visible input \(a\) and closed state
\(x\) there is a unique arrow of \(\mathbb P^\circ\) over \(a\).  By
Theorem~\ref{thm:internal-trace-class-criterion}, this arrow is represented by
the unique feedback witness
\[
h_\Phi(a,x)
\]
satisfying
\[
h_\Phi(a,x)=\omega_U((a,h_\Phi(a,x)),x).
\]
Transporting the target and output components of this unique arrow gives
\[
\delta^{\operatorname{tr}}(a,x)
=
\delta((a,h_\Phi(a,x)),x)
\]
and
\[
\omega^{\operatorname{tr}}(a,x)
=
\omega_B((a,h_\Phi(a,x)),x).
\]

It remains to check that the transition lands in \(P^\circ\).  By the typing
equations for the original Mealy morphism,
\[
p_U\delta((a,h_\Phi(a,x)),x)
=
s_Uh_\Phi(a,x),
\]
while
\[
q_U\delta((a,h_\Phi(a,x)),x)
=
s_U\omega_U((a,h_\Phi(a,x)),x).
\]
Since
\[
h_\Phi(a,x)=\omega_U((a,h_\Phi(a,x)),x),
\]
these two objects agree.  Hence
\[
\delta^{\operatorname{tr}}(a,x)\in P^\circ.
\]
\end{proof}

\section{Restricted trace along the Mealy subcategory}
\label{sec:restricted-trace-along-mealy-subcategory}

The preceding sections define the traced Mealy morphism by applying the total
span trace and then asking whether the resulting span is again input-discrete.
We now record why this gives a genuine partial trace operation on Mealy
morphisms, rather than an ad hoc construction.

Let
\[
\mathcal C:=\mathbf{Span}(\mathbf{Cat}_{\mathrm{int}}(\mathcal E)).
\]
By Proposition~\ref{prop:ambient-span-total-trace}, \(\mathcal C\) is a traced
symmetric monoidal category.  Let
\[
\mathcal S_{\mathrm{Mly}}
\subseteq
\mathcal C
\]
be the symmetric monoidal subcategory of input-discrete labelled spans.

For a morphism
\[
f:
\mathbb A^{\mathrm{op}}\times\mathbb U^{\mathrm{op}}
\longrightarrow
\mathbb B^{\mathrm{op}}\times\mathbb U^{\mathrm{op}}
\]
in \(\mathcal S_{\mathrm{Mly}}\), represented by
\[
\begin{tikzcd}[column sep=large]
\mathbb A^{\mathrm{op}}\times\mathbb U^{\mathrm{op}}
&
\mathbb P
  \arrow[l, "I"']
  \arrow[r, "\Omega"]
&
\mathbb B^{\mathrm{op}}\times\mathbb U^{\mathrm{op}},
\end{tikzcd}
\]
define
\[
\operatorname{Tr}^{\mathbb U}(f)
\]
to be defined precisely when the ambient trace
\[
\operatorname{Tr}^{\mathbb U^{\mathrm{op}}}_{\mathcal C}(f)
\]
belongs again to \(\mathcal S_{\mathrm{Mly}}\).  In that case,
\[
\operatorname{Tr}^{\mathbb U}(f)
:=
\operatorname{Tr}^{\mathbb U^{\mathrm{op}}}_{\mathcal C}(f).
\]
Thus the trace is not an additional operation.  It is the restriction of the
ambient total trace along the inclusion
\[
\mathcal S_{\mathrm{Mly}}\hookrightarrow\mathcal C.
\]

\begin{theorem}[Restricted span trace]
\label{thm:restricted-span-trace}
The operation above defines the partial trace on
\[
\mathcal S_{\mathrm{Mly}}
\]
induced by the ambient total trace on
\[
\mathcal C=
\mathbf{Span}(\mathbf{Cat}_{\mathrm{int}}(\mathcal E)).
\]

For every traced expression built from Mealy morphisms, its value is defined
exactly when the corresponding ambient traced span has input-discrete visible
left leg.  Whenever the relevant traced expressions are trace-class, the trace
equations hold in
\[
\mathcal S_{\mathrm{Mly}}.
\]
Equivalently, \(\mathcal S_{\mathrm{Mly}}\) inherits the trace theory of
\(\mathcal C\) on the full trace-class domain cut out by input-discreteness.
\end{theorem}

\begin{proof}
Every traced expression in
\[
\mathcal S_{\mathrm{Mly}}
\]
has an underlying traced expression in
\[
\mathcal C.
\]
The inclusion
\[
\mathcal S_{\mathrm{Mly}}\hookrightarrow\mathcal C
\]
is faithful and symmetric monoidal.  Hence, whenever a traced expression is
defined in \(\mathcal S_{\mathrm{Mly}}\), its value is exactly the corresponding
ambient traced span.

The trace equations hold in \(\mathcal C\).  Therefore, if a traced equation is
formed from trace-class expressions in \(\mathcal S_{\mathrm{Mly}}\), its two
sides have the same image in \(\mathcal C\).  Faithfulness of the inclusion
then implies equality in \(\mathcal S_{\mathrm{Mly}}\).

It remains to identify the domain of definition.  For a span
\[
\begin{tikzcd}[column sep=large]
\mathbb A^{\mathrm{op}}\times\mathbb U^{\mathrm{op}}
&
\mathbb P
  \arrow[l, "I"']
  \arrow[r, "\Omega"]
&
\mathbb B^{\mathrm{op}}\times\mathbb U^{\mathrm{op}},
\end{tikzcd}
\]
write
\[
I=(I_A,I_U),
\qquad
\Omega=(\Omega_B,\Omega_U).
\]
The ambient trace has apex
\[
\mathbb P^\circ=\operatorname{Eq}(I_U,\Omega_U)
\]
and visible left leg
\[
I_A:\mathbb P^\circ\to\mathbb A^{\mathrm{op}}.
\]
By definition of
\[
\mathcal S_{\mathrm{Mly}},
\]
the traced span belongs to \(\mathcal S_{\mathrm{Mly}}\) exactly when this
visible left leg is input-discrete.  This is precisely the trace-class
condition.
\end{proof}

\section{Simultaneous and staged feedback}
\label{sec:simultaneous-staged-feedback}

The total span trace satisfies the vanishing and nesting laws because
equalizers of products of interfaces may be computed as iterated equalizers.
In the Mealy subcategory, the same statement has a deterministic
interpretation: staged feedback is available exactly when the corresponding
intermediate feedback problems remain uniquely solvable.

Let
\[
\Phi:
\mathbb A\times\mathbb U\times\mathbb V
\rightsquigarrow
\mathbb B\times\mathbb U\times\mathbb V
\]
be a Mealy morphism.  In components, write the hidden output parts as
\[
\omega_U,
\qquad
\omega_V.
\]
A simultaneous closed-loop witness consists of hidden arrows
\[
h_U\in U_1,
\qquad
h_V\in V_1
\]
satisfying
\[
h_U=\omega_U((a,h_U,h_V),x),
\qquad
h_V=\omega_V((a,h_U,h_V),x).
\]

Let
\[
W_\Phi^{U,V}
\]
be the internal object of such simultaneous witnesses, and let
\[
D_\Phi^{U,V}
\]
be the object of visible admissible input/closed-state pairs.  Thus
\[
\Phi
\]
is trace-class over
\[
\mathbb U\times\mathbb V
\]
precisely when
\[
W_\Phi^{U,V}\to D_\Phi^{U,V}
\]
is an isomorphism.

Suppose first that \(\Phi\) is trace-class over \(\mathbb V\), treating
\(\mathbb U\) as part of the visible interface.  Then there is an internal
solution
\[
h_V=h_V(a,h_U,x)
\]
of
\[
h_V=\omega_V((a,h_U,h_V),x),
\]
natural in the visible data \((a,h_U,x)\).  Substitution gives a traced Mealy
morphism
\[
\operatorname{Tr}^{\mathbb V}(\Phi):
\mathbb A\times\mathbb U
\rightsquigarrow
\mathbb B\times\mathbb U.
\]
Its remaining \(\mathbb U\)-feedback equation is
\[
h_U=
\omega_U((a,h_U,h_V(a,h_U,x)),x).
\]

\begin{theorem}[Staged feedback theorem]
\label{thm:staged-feedback}
Let
\[
\Phi:
\mathbb A\times\mathbb U\times\mathbb V
\rightsquigarrow
\mathbb B\times\mathbb U\times\mathbb V
\]
be a Mealy morphism.

If
\[
\operatorname{Tr}^{\mathbb V}(\Phi)
\]
is defined and
\[
\operatorname{Tr}^{\mathbb U}
\bigl(\operatorname{Tr}^{\mathbb V}(\Phi)\bigr)
\]
is defined, then the simultaneous trace over
\[
\mathbb U\times\mathbb V
\]
is defined and there is a canonical equality
\[
\operatorname{Tr}^{\mathbb U}
\bigl(\operatorname{Tr}^{\mathbb V}(\Phi)\bigr)
=
\operatorname{Tr}^{\mathbb U\times\mathbb V}(\Phi).
\]

On components, the staged solution
\[
h_V=h_V(a,h_U,x)
\]
and then
\[
h_U=h_U(a,x)
\]
determine the simultaneous solution
\[
(h_U(a,x),h_V(a,h_U(a,x),x)).
\]
Conversely, whenever the simultaneous solution factors through a unique
intermediate \(\mathbb V\)-solution natural in \((a,h_U,x)\), the staged trace
is defined and agrees with the simultaneous trace.
\end{theorem}

\begin{proof}
In the ambient span category, the equality
\[
\operatorname{Tr}^{\mathbb U}
\bigl(\operatorname{Tr}^{\mathbb V}(-)\bigr)
=
\operatorname{Tr}^{\mathbb U\times\mathbb V}(-)
\]
is the vanishing-II law of the total trace.

The apex of the simultaneous trace is the equalizer imposing both hidden
matching equations
\[
I_U=\Omega_U,
\qquad
I_V=\Omega_V.
\]
Equivalently, it is the iterated equalizer obtained by first imposing
\[
I_V=\Omega_V
\]
and then imposing
\[
I_U=\Omega_U
\]
on the result.  Since finite limits in
\[
\mathbf{Cat}_{\mathrm{int}}(\mathcal E)
\]
are computed componentwise, these two constructions give canonically
isomorphic internal categories.

If the staged trace is defined, then after the first equalizer the remaining
visible left leg is input-discrete, and after the second equalizer the final
visible left leg is input-discrete.  Hence the staged trace is a morphism of
\[
\mathcal S_{\mathrm{Mly}}.
\]
The ambient vanishing-II isomorphism identifies this staged traced span with
the simultaneous traced span.  Therefore the simultaneous traced span is also
input-discrete over the final visible input, and the two Mealy morphisms agree.

In component notation, the first input-discreteness condition says that for
each visible tuple \((a,h_U,x)\) there is a unique
\[
h_V=h_V(a,h_U,x)
\]
solving
\[
h_V=\omega_V((a,h_U,h_V),x).
\]
The second says that for each \((a,x)\) there is a unique
\[
h_U=h_U(a,x)
\]
solving
\[
h_U=\omega_U((a,h_U,h_V(a,h_U,x)),x).
\]
Substitution gives the unique simultaneous pair
\[
(h_U(a,x),h_V(a,h_U(a,x),x)),
\]
which satisfies both hidden feedback equations.  The visible transition and
output obtained by substituting this pair are exactly those obtained by the
ambient simultaneous trace.
\end{proof}

\section{The partial trace theorem}
\label{sec:partial-trace-theorem}

\begin{theorem}[Partial trace on Mealy morphisms]
\label{thm:mly-partially-traced}
The symmetric monoidal category
\[
\mathsf{Mly}(\mathcal E)
\]
carries the canonical partial trace obtained by restricting the total trace of
\[
\mathbf{Span}(\mathbf{Cat}_{\mathrm{int}}(\mathcal E))
\]
to the input-discrete span representation of Mealy morphisms.

Explicitly, for
\[
\Phi:
\mathbb A\times\mathbb U
\rightsquigarrow
\mathbb B\times\mathbb U,
\]
the trace
\[
\operatorname{Tr}^{\mathbb U}(\Phi):
\mathbb A\rightsquigarrow\mathbb B
\]
is defined precisely when \(\Phi\) is trace-class over \(\mathbb U\), i.e.
when the closed-loop visible input projection
\[
I_A:\mathbb P^\circ\to\mathbb A^{\mathrm{op}}
\]
is input-discrete.  Equivalently, the internal feedback projection
\[
\pi_\Phi:W_\Phi\to D_\Phi
\]
is an isomorphism.

When defined, the trace is represented by the ambient span trace
\[
\begin{tikzcd}[column sep=large]
\mathbb A^{\mathrm{op}}
&
\mathbb P^\circ
  \arrow[l, "I_A"']
  \arrow[r, "\Omega_B"]
&
\mathbb B^{\mathrm{op}} .
\end{tikzcd}
\]
The trace equations are inherited from the total trace on spans on their
trace-class domains.  In particular, staged feedback agrees with simultaneous
feedback whenever the staged feedbacks are trace-class.
\end{theorem}

\begin{proof}
By Theorem~\ref{thm:mealy-as-input-discrete-spans}, Mealy morphisms are
faithfully represented as the symmetric monoidal subcategory
\[
\mathcal S_{\mathrm{Mly}}
\subseteq
\mathbf{Span}(\mathbf{Cat}_{\mathrm{int}}(\mathcal E))
\]
of labelled spans whose left leg is input-discrete.

The ambient span category is traced by
Proposition~\ref{prop:ambient-span-total-trace}.  For a Mealy morphism
\[
\Phi:
\mathbb A\times\mathbb U
\rightsquigarrow
\mathbb B\times\mathbb U,
\]
the ambient trace of its labelled span has apex
\[
\mathbb P^\circ=\operatorname{Eq}(I_U,\Omega_U)
\]
and visible left leg
\[
I_A:\mathbb P^\circ\to\mathbb A^{\mathrm{op}}.
\]
The traced span lies again in
\[
\mathcal S_{\mathrm{Mly}}
\]
exactly when this left leg is input-discrete.  This is the trace-class
condition.

The equivalence with
\[
\pi_\Phi:W_\Phi\to D_\Phi
\]
being an isomorphism is Theorem~\ref{thm:internal-trace-class-criterion}.
The component formula is Proposition~\ref{prop:component-formula-partial-trace}.

Finally, every trace equation for trace-class Mealy morphisms is the image of
the corresponding trace equation in the ambient traced span category.  Since
the inclusion of
\[
\mathcal S_{\mathrm{Mly}}
\]
into the ambient span category is faithful, the equality holds in
\[
\mathsf{Mly}(\mathcal E).
\]
The staged-feedback statement is Theorem~\ref{thm:staged-feedback}.
\end{proof}

\section{Interface privacy}
\label{sec:interface-privacy}

The trace operation has a simple operational interpretation.  A Mealy morphism
\[
\Phi:
\mathbb A\times\mathbb U
\rightsquigarrow
\mathbb B\times\mathbb U
\]
has a visible interface
\[
\mathbb A\rightsquigarrow\mathbb B
\]
and a hidden interface
\[
\mathbb U\rightsquigarrow\mathbb U.
\]
Tracing over \(\mathbb U\) hides the internal \(\mathbb U\)-conversation by
forcing the hidden output to match the hidden input.

The ambient span trace remembers all hidden witnesses satisfying this matching
condition.  The Mealy trace exists exactly when these hidden witnesses are
unique over each visible input and closed state.  Thus trace-class feedback is
the condition under which the hidden interface can be eliminated without
leaving the deterministic Mealy world.

In component form, the hidden witness is the solution of
\[
h=\omega_U((a,h),x).
\]
When this solution is unique, the visible behaviour is obtained by substituting
the solution back into the original Mealy transition and output:
\[
\delta^{\operatorname{tr}}(a,x)
=
\delta((a,h_\Phi(a,x)),x),
\qquad
\omega^{\operatorname{tr}}(a,x)
=
\omega_B((a,h_\Phi(a,x)),x).
\]

\section{Feedback and relative execution semantics}
\label{sec:feedback-relative-execution-semantics}

We finally record how the partial trace interacts with the relative execution
semantics developed earlier.

Let
\[
\Phi:
\mathbb A\times\mathbb U
\rightsquigarrow
\mathbb B\times\mathbb U
\]
be trace-class over \(\mathbb U\), and let
\[
Y
\]
be a right \(\mathbb B\)-action.  The traced Mealy morphism
\[
\operatorname{Tr}^{\mathbb U}(\Phi):
\mathbb A\rightsquigarrow\mathbb B
\]
therefore determines a relative program category
\[
\mathsf{Prog}_{\operatorname{Tr}^{\mathbb U}(\Phi),Y}.
\]

The states of this program category are closed relative states
\[
(x,y)
\]
with
\[
x\in P^\circ,
\qquad
q_B(x)=r(y).
\]
For an admissible visible input
\[
a:u\to p_A(x),
\]
the transition is computed by first solving the feedback equation
\[
h_\Phi(a,x)=\omega_U((a,h_\Phi(a,x)),x),
\]
and then applying the original relative transition:
\[
(a,x,y)
\longmapsto
\left(
\delta((a,h_\Phi(a,x)),x),
\,
\sigma(\omega_B((a,h_\Phi(a,x)),x),y)
\right).
\]

Thus relative execution after feedback is obtained by closing the hidden
interface first and then running the resulting visible Mealy morphism.

\begin{proposition}[Relative execution after feedback]
\label{prop:relative-execution-after-feedback}
Let
\[
\Phi:
\mathbb A\times\mathbb U
\rightsquigarrow
\mathbb B\times\mathbb U
\]
be trace-class over \(\mathbb U\), and let \(Y\) be a right
\(\mathbb B\)-action.  Then the relative program category of
\[
\operatorname{Tr}^{\mathbb U}(\Phi)
\]
has object object
\[
P^\circ\times_{q_B,B_0,r}Y
\]
and arrow object
\[
A_1\times_{t_A,A_0,p_A}
\left(P^\circ\times_{q_B,B_0,r}Y\right).
\]
Its source, target, identities, and composition are those induced by the
transition and output maps of
\[
\operatorname{Tr}^{\mathbb U}(\Phi).
\]
Equivalently, its one-step execution is computed by substituting the unique
feedback solution \(h_\Phi(a,x)\) into the original Mealy transition.
\end{proposition}

\begin{proof}
This is the relative program-category construction applied to the traced Mealy
morphism.  The carrier of
\[
\operatorname{Tr}^{\mathbb U}(\Phi)
\]
is
\[
A_0\xleftarrow{p_A}P^\circ\xrightarrow{q_B}B_0,
\]
and its transition and output maps are given by
Proposition~\ref{prop:component-formula-partial-trace}.  Substituting these
maps into the definition of the relative action yields the displayed object
and arrow objects and the stated one-step transition formula.

The emitted visible output arrow is
\[
\omega^{\operatorname{tr}}(a,x)
=
\omega_B((a,h_\Phi(a,x)),x),
\]
and the downstream state is obtained by applying the \(Y\)-action to this
arrow:
\[
\sigma(\omega^{\operatorname{tr}}(a,x),y).
\]
\end{proof}

\section{Conceptual summary}
\label{sec:partial-trace-conceptual-summary}

The construction of this chapter can be summarized in two equivalent ways.

First, categorically:
\[
\boxed{
\mathsf{Mly}(\mathcal E)
\text{ is a symmetric monoidal subcategory of }
\mathbf{Span}(\mathbf{Cat}_{\mathrm{int}}(\mathcal E)),
}
\]
and the latter has a total trace.  The Mealy trace is the restriction of this
ambient trace to the morphisms whose traced spans remain input-discrete.

Second, operationally:
\[
\boxed{
\operatorname{Tr}^{\mathbb U}(\Phi)
\text{ is defined exactly when the feedback equation }
h=\omega_U((a,h),x)
\text{ has a unique solution.}
}
\]
When this holds, the traced transition and output are obtained by substituting
the unique feedback solution into the original Mealy transition:
\[
\delta^{\operatorname{tr}}(a,x)
=
\delta((a,h_\Phi(a,x)),x),
\qquad
\omega^{\operatorname{tr}}(a,x)
=
\omega_B((a,h_\Phi(a,x)),x).
\]

For multiple hidden interfaces, staged feedback agrees with simultaneous
feedback on the trace-class staged domain:
\[
\operatorname{Tr}^{\mathbb U}
\bigl(\operatorname{Tr}^{\mathbb V}(\Phi)\bigr)
=
\operatorname{Tr}^{\mathbb U\times\mathbb V}(\Phi)
\]
whenever the staged traces are defined.  This is the deterministic Mealy form
of the vanishing law for the compact-closed trace of spans.

Thus the partial trace is not an additional axiom on Mealy morphisms.  It is
the canonical feedback operation inherited from spans, restricted by the
requirement that closed-loop execution remain input-discrete.